\documentclass[12pt]{article}
\usepackage{amsmath, amssymb, amsfonts, amsthm, geometry, setspace, tcolorbox}
\geometry{margin=1in}
\setstretch{1.2}
\setlength{\parindent}{0pt}

\tcbset{colback=gray!5, colframe=black!70, boxrule=0.6pt, arc=2pt, 
left=6pt, right=6pt, top=4pt, bottom=4pt}

\newtheorem{theorem}{Theorem}
\newtheorem{definition}{Definition}

\title{\textbf{Envelope Theorem and Duality}\\[7pt]
\large ECON 320 - Introduction to Mathematical Economics}
\author{Muhebullah Karimzada}
\date{Fall 2025}

\begin{document}
\maketitle



\vspace{1em}

\begin{tcolorbox}
\textbf{Learning Objectives}
\begin{itemize}
\item Understand the envelope theorem for constrained and unconstrained problems.
\item Apply the theorem in consumer, producer, and profit-maximization settings.
\item Develop intuition for duality in economic optimization.
\item Learn the relationship between Lagrange multipliers, shadow prices, and value functions.
\end{itemize}
\end{tcolorbox}

%-------------------------------------------------------------
\section*{1. Introduction}

Many economic optimization problems depend on an external parameter (price, income, technology). 
We frequently want to know how the \textit{optimal value} of the objective changes when the parameter changes, without re-solving the optimization problem completely.

The \textbf{Envelope Theorem} provides an elegant method to compute this derivative.  
The related concept of \textbf{Duality} expresses an optimization problem in an equivalent but economically insightful form.

%-------------------------------------------------------------
\section*{2. The Optimization Setup}

Consider the constrained maximization problem:
\[
\max_{x_1,\ldots,x_n} f(x_1,\ldots,x_n,a)
\]
subject to
\[
g(x_1,\ldots,x_n,a)=0,
\]
where $a$ is a parameter.

Let $x_i^*(a)$ denote the optimal choices, and let the \textbf{value function} be
\[
V(a)=f(x_1^*(a), \ldots, x_n^*(a), a).
\]

%-------------------------------------------------------------
\section*{3. The Envelope Theorem}

\begin{theorem}[Envelope Theorem — Constrained Case]
For the constrained maximization problem above, the derivative of the value function is:
\[
\boxed{
\frac{dV}{da} = 
\frac{\partial f}{\partial a}(x^*,a)
-
\lambda \frac{\partial g}{\partial a}(x^*,a)
}
\]
where $\lambda$ is the optimal Lagrange multiplier.
\end{theorem}

\textbf{Key idea:} At the optimum, the indirect effects through $dx_i^*/da$ cancel out by the first-order conditions.

\subsection*{Unconstrained Version}
For:
\[
\max_x f(x,a),
\quad f_x(x^*,a)=0,
\]
\[
\boxed{
\frac{dV}{da}
=
\frac{\partial f}{\partial a}(x^*(a),a).
}
\]

%-------------------------------------------------------------
\section*{4. Economic Interpretation of the Envelope Theorem}

\begin{tcolorbox}
\textbf{Interpretation:} 
Only the \textit{direct effect} of the parameter on the objective and constraint matters.  
All indirect effects through optimal choices vanish.
\end{tcolorbox}

Examples:
\begin{itemize}
\item In consumer theory: $\displaystyle \frac{dU^*}{dM}=\lambda$ = marginal utility of income.
\item In profit maximization: $\displaystyle \frac{\partial \pi^*}{\partial w}=-L^*$ = negative of optimal labor usage.
\item In cost minimization: $\displaystyle \frac{dC}{dQ}=\lambda$ = marginal cost.
\end{itemize}

%-------------------------------------------------------------
\section*{5. Applications}

\subsection*{5.1 Consumer Utility Maximization}
\[
\max_{x,y} U(x,y) \quad \text{s.t. } p_x x + p_y y = M.
\]
Lagrangian:
\[
\mathcal{L}=U(x,y)+\lambda(M-p_x x-p_y y).
\]
Envelope Theorem:
\[
\frac{dU^*}{dM}=\lambda.
\]
Thus $\lambda$ is the \textbf{marginal utility of income}.

%-------------------------------------------------------------
\subsection*{5.2 Profit Maximization}
\[
\pi=\max_{K,L}[pF(K,L)-wL-rK].
\]
Envelope Theorem gives:
\[
\frac{\partial \pi}{\partial w}=-L^*,\quad 
\frac{\partial \pi}{\partial r}=-K^*,\quad
\frac{\partial \pi}{\partial p}=F(K^*,L^*).
\]
Thus:
\begin{itemize}
\item Profit falls with each input price at the rate of input usage.
\item Profit rises with output price at the rate of output.
\end{itemize}

%-------------------------------------------------------------
\subsection*{5.3 Cost Minimization}
\[
C(w,r,Q)=\min_{K,L}[wL+rK]\quad\text{s.t.}\quad F(K,L)=Q.
\]
Envelope Theorem:
\[
\frac{\partial C}{\partial Q}=\lambda=\text{marginal cost}.
\]

%-------------------------------------------------------------
\section*{6. Duality}

\begin{definition}
\textbf{Duality} refers to the correspondence between an optimization problem (the \emph{primal}) and another optimization problem (the \emph{dual}) that characterizes the same economic behavior in a different form.
\end{definition}

%-------------------------------------------------------------
\subsection*{6.1 Consumer Duality}

\begin{center}
\begin{tabular}{c|c|c}
\textbf{Problem} & \textbf{Objective} & \textbf{Constraint} \\ \hline
Utility Maximization & Maximize $U(x,y)$ & $p_x x + p_y y = M$ \\
Expenditure Minimization & Minimize $p_x x + p_y y$ & $U(x,y)=\bar{U}$
\end{tabular}
\end{center}

Both yield the same tangency condition:
\[
MRS=\frac{p_x}{p_y}.
\]

%-------------------------------------------------------------
\subsection*{6.2 Producer Duality}
\begin{center}
\begin{tabular}{c|c|c}
\textbf{Problem} & \textbf{Objective} & \textbf{Constraint} \\ \hline
Cost Minimization & $\min\;C=wL+rK$ & $F(K,L)=Q$ \\
Profit Maximization & $\max\;\pi=pF(K,L)-wL-rK$ & None
\end{tabular}
\end{center}

Key identity:
\[
\pi(p,w,r)=pQ-C(w,r,Q).
\]

%-------------------------------------------------------------
\section*{7. Shephard’s Lemma}

For the cost function:
\[
C(w,r,Q)=\min_{K,L}\{wL+rK: F(K,L)=Q\},
\]
\[
\boxed{
\frac{\partial C}{\partial w}=L^*, \qquad 
\frac{\partial C}{\partial r}=K^*
}
\]
That is, the derivative of the cost function with respect to input prices returns the \textbf{conditional factor demands}.

This is a direct application of the Envelope Theorem.

%-------------------------------------------------------------
\section*{8. Graphical Insight}

\begin{itemize}
\item The value function forms an \textbf{envelope} of the objective curves for different $x$.
\item Each curve touches the envelope at the optimum.
\item The envelope's slope gives the derivative of optimal value with respect to a parameter.
\end{itemize}

%-------------------------------------------------------------
\section*{9. Summary of Key Results}

\begin{tcolorbox}
\[
\frac{dV}{da} = \frac{\partial f}{\partial a} 
- \lambda\frac{\partial g}{\partial a}
\]
\[
\lambda = \text{shadow value of the constraint}
\]
\[
\frac{\partial C}{\partial Q}=\text{marginal cost}
\]
\[
\frac{\partial \pi}{\partial w}=-L^*, \quad 
\frac{\partial \pi}{\partial r}=-K^*
\]
\[
\frac{\partial C}{\partial w}=L^*, \qquad \frac{\partial C}{\partial r}=K^*
\]
\end{tcolorbox}

Envelope Theorem + Duality form one of the most powerful toolkits for modern microeconomic analysis.

\end{document}
